\subsection{Canalización de Ingesta Adaptativa y Fusión Temprana Audio-Visual}
\label{subsec:pipeline_ingesta_fasion_temprana}

El procesamiento de secuencias de vídeo continuas de larga duración requiere segmentar el metraje en unidades semánticas discretas que preserven la coherencia visual y auditiva, respetando estrictamente las cuotas de inferencia de la API fundacional ($100\text{ RPM}$, $30.000\text{ TPM}$ y $1.000\text{ RPD}$) \cite{gemini_emb2_docs}.

\subsubsection{Segmentación Temporal Adaptativa mediante Similitud Estructural (SSIM)}
En lugar de emplear ventanas temporales fijas que fragmentan arbitrariamente eventos continuos o introducen mezclas incoherentes de planos, el sistema implementa una segmentación adaptativa de una sola pasada (\textit{single-pass}). 

Para cada par de fotogramas secuenciales en escala de grises $I_{t-1}, I_t$, se evalúa el Índice de Similitud Estructural (\textit{Structural Similarity Index Measure}, SSIM):
\begin{equation}
    \text{SSIM}(I_{t-1}, I_t) = \frac{(2\mu_{t-1}\mu_t + C_1)(2\sigma_{t-1,t} + C_2)}{(\mu_{t-1}^2 + \mu_t^2 + C_1)(\sigma_{t-1}^2 + \sigma_t^2 + C_2)}
\end{equation}
donde $\mu$ y $\sigma^2$ representan la media y varianza local de luminancia, $\sigma_{t-1,t}$ es la covarianza cruzada, y $C_1 = (K_1 L)^2, C_2 = (K_2 L)^2$ son constantes estabilizadoras ante regiones oscuras ($L=255, K_1=0,01, K_2=0,03$).

Un corte de plano se declara formalmente en el instante $t$ si se satisface la condición de transición abrupta o se alcanza la cota superior temporal de coherencia:
\begin{equation}
    \text{Corte}(t) = \left( \text{SSIM}(I_{t-1}, I_t) < \theta_{\text{SSIM}} \land \Delta t_{\text{actual}} \ge T_{\min} \right) \lor \left( \Delta t_{\text{actual}} \ge T_{\max} \right)
\end{equation}
donde $\theta_{\text{SSIM}} = 0,65$, $T_{\min} = 1,0\text{ s}$ y $T_{\max} = 12,0\text{ s}$. Para cada plano detectado, se extraen tres fotogramas clave correspondientes al instante inicial ($t_{\text{start}}$), medio ($t_{\text{mid}}$) y final ($t_{\text{end}}$). Cada fotograma es escalado a un ancho uniforme de 512 píxeles mediante interpolación bilineal y codificado en formato JPEG (calidad 80\%), reduciendo sustancialmente el tamaño de transferencia sin degradar la geometría espacial.

\subsubsection{Extracción y Alineación Forzada de la Locución Auditiva}
De forma paralela a la segmentación visual, la pista de audio completa del contenedor de vídeo es desacoplada mediante \texttt{ffmpeg}, remuestreada a $16\text{ kHz}$ en formato PCM monocanal y procesada con el modelo de reconocimiento automático del habla OpenAI Whisper (\textit{checkpoint} \texttt{base}) \cite{whisperx_silero}.

Whisper produce un flujo de segmentos transcritos $\mathcal{A} = \{A_1, A_2, \dots, A_M\}$, donde cada segmento $A_k = (t_{a,k}^{\text{start}}, t_{a,k}^{\text{end}}, w_k)$ incluye marcas temporales de alta resolución y el texto fonético correspondiente $w_k$. La alineación forzada entre los planos visuales $C_i = [t_i^{\text{start}}, t_i^{\text{end}}]$ y las locuciones auditivas se establece evaluando el solape temporal efectivo:
\begin{equation}
    \text{Solape}(C_i, A_k) = \max\left(0, \min(t_i^{\text{end}}, t_{a,k}^{\text{end}}) - \max(t_i^{\text{start}}, t_{a,k}^{\text{start}})\right)
\end{equation}

Un segmento de locución se vincula al fragmento visual $C_i$ si cumple el umbral mínimo de solape $\text{Solape}(C_i, A_k) \ge \Delta_{\min} = 0,2\text{ s}$. La transcripción global $\tau_i$ asociada al plano visual resulta de la concatenación cronológica de todos los segmentos vinculados:
\begin{equation}
    \tau_i = \bigoplus_{k: \text{Solape}(C_i, A_k) \ge \Delta_{\min}} w_k
\end{equation}

\subsubsection{Fusión Temprana Omnimodal en Solicitud Única}
Para cumplir con los límites de cuota de la capa gratuita sin renunciar a la riqueza perceptual de ambas modalidades, la inyección de la información visual y auditiva se efectúa mediante fusión temprana (\textit{early fusion}). 

Los tres fotogramas clave codificados en JPEG se transforman en objetos binarios estructurados de tipo imagen (\texttt{types.Part.from\_bytes}), mientras que la locución alineada $\tau_i$ se incorpora como una parte textual explícita (\texttt{types.Part.from\_text}) dentro del mismo contenedor unificado \texttt{types.Content}. La petición a la API de \textit{Gemini Embedding 2} se formula en un único paquete atómico:
\begin{equation}
    \text{Payload}_i = \left\{ \text{parts}: \left[ \text{JPEG}(f_{i,1}), \text{JPEG}(f_{i,2}), \text{JPEG}(f_{i,3}), \text{Text}(\tau_i) \right] \right\}
\end{equation}

Este diseño garantiza que cada fragmento audiovisual consume exactamente una única llamada a la API de vectorización, preservando una estricta invarianza en el coste de cuota ($1\text{ chunk} = 1\text{ request}$) con independencia de la duración o complejidad del segmento.
